\section{数据集、预处理与训练}
\label{sec:experiment}
\subsection{数据集}
项目组合使用通用表情样本与车载近红外样本。RAF-DB 提供七类自然场景表情监督，KMU-FED 提供真实车辆座舱近红外图像，PPB-Emo 中央红外视频补充中性样本，FER+ 的投票标签用于筛选一致性较高的补充图像。各数据集分工见表\ref{tab:datasets}。

\begin{table}[htbp]
\centering
\caption{实际使用的数据与训练用途}
\label{tab:datasets}
\begin{tabularx}{\linewidth}{@{}lXX@{}}
\toprule
数据集 & 本项目使用规模与类别 & 用途\\
\midrule
RAF-DB & 12\,271 张训练、3\,068 张验证；七类 & 通用表情特征与验证\\
KMU-FED & 1\,100 张近红外人脸，12 名参与者；六类 & 座舱域适配与匿名对照\\
PPB-Emo & 对齐清单 185 帧，采用其中 183 帧中性样本 & 近红外中性补充\\
FER2013/FER+ & 35\,887 行原始像素与投票标签，经筛选使用 & 表情与中性标签补充\\
\bottomrule
\end{tabularx}
\end{table}

RAF-DB 的基本表情子集覆盖惊讶、恐惧、厌恶、高兴、悲伤、愤怒和中性\cite{li2017rafdb}。KMU-FED 的近红外采集方式与车辆使用场景相符\cite{kmu2018}。PPB-Emo 是驾驶任务中的多模态数据集，本项目使用中央红外面部视频与对应中性标注\cite{ppbemo}。FER+ 在原有图像上补充多名标注者投票\cite{barsoum2016ferplus}，数据读取时要求严格多数和至少五票，并统一映射到七类标签。

\subsection{预处理与划分}
输入先转为单通道灰度图，人脸经五点相似变换对齐。KMU-FED 与 PPB-Emo 的对齐图以 $256\times256$ 保存，实际分类输入为匿名后重新对齐的 $100\times100$ 灰度脸。归一化统计取自训练集，验证不使用随机增强；质量配置按模型契约执行 $260\times260$ 浮点插值和 ImageNet 归一化。

RAF-DB 使用官方训练与验证划分。KMU-FED 按序列保留尾帧，形成 880 张训练样本和 220 张回归样本，评价范围为既定人员与序列。PPB-Emo 按参与者划分为 114 张训练、34 张验证和 35 张测试图像，三部分人员不重叠。划分完成后再生成匿名缓存，保持源图、标签与数据用途的对应关系。

\subsection{表情模型训练与校准}
训练分为通用基线、多域联合和匿名域微调三个步骤。基线使用交叉熵学习表情类别；联合训练通过加权采样引入 KMU-FED、FER+ 与 PPB-Emo；匿名域阶段直接使用运行链路生成的匿名灰度缓存，让分类输入与应用中的预处理保持一致。

\begin{table}[htbp]
\centering
\caption{训练器配置与各阶段任务}
\label{tab:training}
\begin{tabularx}{\linewidth}{@{}lXX@{}}
\toprule
阶段 & 主要配置 & 训练或选择方式\\
\midrule
通用基线 & AdamW，学习率 $10^{-3}$，批量 128，默认至多 20 轮 & 按验证准确率保留权重\\
多域联合 & AdamW，学习率 $10^{-4}$，批量 128，默认 12 轮 & 分域加权采样与联合验证\\
匿名域微调 & 每类最多 650 张 RAF 训练图起步，加入车载样本 & 标签监督与 0.3 权重一致性损失\\
质量配置 & 批量 16，骨干/分类头学习率 $3\times10^{-5}/2\times10^{-4}$ & 微调后段特征与分类头\\
\bottomrule
\end{tabularx}
\end{table}

匿名域训练对同一匿名图构造轻微亮度与平移视图，既要求分类结果符合真实标签，也要求不同视图的输出接近。验证时分别检查 RAF-DB、KMU-FED 和 PPB-Emo，按各域表现选取候选，而非只依据训练损失下降替换运行权重。

质量配置在匿名输入上微调特征与分类头，并对原视图和左右翻转视图的概率求平均。温度校准调整概率分布的尺度，类别由校准后的最大概率给出；模型契约保存类别顺序、预处理、温度和文件校验值。当前导出契约记录选中第 2 轮候选，其验证结果在第五章分别列出。

\subsection{运行环境}
系统基于 Python、PyTorch、OpenCV、PyAV、MediaPipe 与 PySide6，GPU 推理使用 CUDA。视频记录包含帧号、处理状态、匿名输入来源、表情概率及计时信息。性能实验使用 RTX 3060 Laptop GPU，将模型处理时间与视频解码、界面显示、输出编码分别记录。

运行资源按检测、匿名化与表情识别组织。模型加载后复用实例和缓存，连续视频适当复用检测与跟踪结果；图片页缓存已完成样本的结果，切换样本时直接读取。桌面显示与分析线程分离，使重新选择图片、切换页面和停止视频保持响应。
